\section{Introducción}

Los proyectos de minería de datos requieren más que ejecutar algoritmos de aprendizaje automático: necesitan comprender a profundidad la naturaleza de las fuentes, auditar rigurosamente su calidad, preparar los microdatos mediante reglas explícitas y reproducibles, y comunicar con honestidad las limitaciones epistemológicas de cada resultado. La metodología estándar CRISP-DM (\textit{Cross-Industry Standard Process for Data Mining}) y el marco KDD (\textit{Knowledge Discovery in Databases}) conciben un ciclo iterativo estructurado que vincula estrechamente la comprensión del negocio o dominio, la exploración y preparación de los datos, el modelado, la evaluación y el despliegue operativo \parencite{chapman2000crisp, fayyad1996data}. En el análisis de microdatos sociodemográficos procedentes de encuestas por muestreo probabilístico, esta rigurosidad es indispensable: códigos categóricos, valores especiales de no respuesta y filtros o saltos condicionales del cuestionario pueden inducir a conclusiones gravemente distorsionadas si se procesan como atributos numéricos convencionales o si se aplican recetas simplistas de imputación masiva.

El presente proyecto aborda el procesamiento integral de \texttt{data/persona.csv}, archivo de microdatos de personas de la Encuesta de Hogares de Bolivia (EH2025). La estructura local observada consta de 39.497 registros individuales y 275 variables. Se cuenta con un diccionario de metadatos de referencia basado en el archivo F27 del catálogo oficial ANDA del Instituto Nacional de Estadística (INE); no obstante, el catálogo oficial declara 39.485 casos y 276 variables \parencite{ineeh2025persona}. Esta discrepancia (12 registros adicionales y variaciones en la denominación de campos) prohíbe asumir apriorísticamente que la copia de trabajo sea idéntica a la publicación oficial, exigiendo un principio de conservación y auditoría transparente sobre la fuente efectivamente analizada.

\subsection{Problema y pregunta de trabajo}
En el ámbito de la ciencia de datos aplicada a políticas públicas, el problema de limpieza no consiste simplemente en sustituir celdas vacías o eliminar observaciones atípicas. Consiste en generar una base analítica reproducible sin destruir información cuya interpretación responde al diseño muestral de la encuesta. En este contexto:
\begin{itemize}[leftmargin=0.5in]
  \item Una celda con token de ausencia (\texttt{NA}) no equivale necesariamente a un error del encuestador o una omisión del informante: con frecuencia representa una persona que legítimamente no forma parte del universo de la pregunta (por ejemplo, preguntas de empleo para menores de 7 años o de maternidad para varones).
  \item Un valor extremo en distribuciones asimétricas de ingreso o edad puede ser perfectamente plausible y no constituye un error que deba suprimirse mediante métodos de truncamiento ciego como el rango intercuartílico (IQR).
  \item Un código numérico discreto (como el departamento o la condición de actividad) no posee propiedades métricas y no admite cálculos de media o desviación estándar sin provocar un error categorial.
\end{itemize}

Frente a este escenario, la pregunta central que guía el proyecto es: \textit{¿Cómo estructurar un pipeline determinista y auditable que limpie y valide técnicamente el dataset \texttt{persona.csv}, preservando la inmutabilidad de la fuente cruda, respetando la estructura de universos muestrales y desplegando un sistema analítico interactivo en Flask con persistencia atómica en JSON local?}

\subsection{Objetivos}
El objetivo general es diseñar, ejecutar y documentar un sistema verificable de preparación, auditoría de calidad y visualización analítica de \texttt{persona.csv}, garantizando la inmutabilidad del archivo original, aplicando transformaciones semánticas trazables y estructurando 11 vistas temáticas especializadas.

Los objetivos específicos son:
\begin{enumerate}[leftmargin=0.5in]
  \item Caracterizar la estructura técnica, la integridad criptográfica (SHA-256) y la taxonomía de ausencias del archivo \texttt{data/persona.csv}.
  \item Validar la unicidad de la clave primaria compuesta \texttt{(folio, nro)} y formalizar el contrato de datos para las 275 columnas.
  \item Implementar y registrar en bitácora transaccional únicamente reglas de limpieza con fundamento físico o lingüístico verificable (reglas S-01 y S-02), preservando el esquema maestro en formato Parquet y CSV.
  \item Construir 11 vistas analíticas temáticas orientadas a subpoblaciones elegibles según los saltos del cuestionario oficial del INE.
  \item Desarrollar una aplicación web interactiva en Flask con arquitectura desacoplada, persistencia atómica en JSON local (sin base de datos relacional pesada) y un módulo geográfico interactivo de los 9 departamentos de Bolivia con datos observados exactos.
\end{enumerate}

\subsection{Alcance y delimitación}
El análisis se concentra exclusivamente en el microdato individual de \texttt{persona.csv}; no realiza fusiones con las bases complementarias de vivienda, gastos o equipamiento del hogar. Las 11 vistas temáticas generadas son proyecciones analíticas de la versión maestra y no descartan variables del esquema global. El alcance abarca preparación, auditoría, perfiles de calidad y visualización descriptiva. No se entrenan modelos de aprendizaje automático predictivo ni se efectúan inferencias poblacionales no ponderadas, en estricto cumplimiento de los supuestos del diseño muestral complejo.
